\par\noindent\textbf{(1)} 引入约化质量 \(\mu=m_1m_2/(m_1+m_2)\)、质心速度 \(\boldsymbol V=(m_1\boldsymbol v_1+m_2\boldsymbol v_2)/(m_1+m_2)\)。动能分解为 \(T=\tfrac12(m_1+m_2)V^2+\tfrac12\mu v_{\mathrm{rel}}^2\)。硬球碰撞只反转相对速度的接触法向分量，切向分量不变。
\par\noindent\textbf{(2)} 接触点的法向与入射速度夹角 \(\alpha\) 满足 \(b=R\sin\alpha\)。镜面反射使相对轨迹的偏转角 \(\chi=\pi-2\alpha=2\arccos(b/R)\)。因此 \(b=R\cos(\chi/2)\)，中心碰撞为 \(\chi=\pi\)，擦边碰撞为 \(\chi=0\)。
\par\noindent\textbf{(3)} 轴对称经典公式 \(d\sigma/d\Omega=(b/\sin\chi)|db/d\chi|\)。代入 \(db/d\chi=-(R/2)\sin(\chi/2)\)，以及 \(\sin\chi=2\sin(\chi/2)\cos(\chi/2)\)，得到 \(d\sigma/d\Omega=R^2/4\)，角分布在中心系各向同性。
\par\noindent\textbf{(4)} 积分所有散射立体角给 \(\sigma=4\pi(R^2/4)=\pi R^2\)，恰是几何截面。若目标球初始静止，实验室散射角与中心系 \(\chi\) 的映射取决于 \(m_1/m_2\)，不能无条件把质心系的常数微分截面直接用于实验室角分布。
